% VeriMem --- ACL-format report.
%
% Compiling: needs acl.sty and acl_natbib.bst beside this file (both are in the
% repo, fetched from github.com/acl-org/acl-style-files). Builds clean with
%     tectonic -X compile verimem.tex
% or on Overleaf from the ACL template. Do not add \bibliographystyle: acl.sty
% already issues it, and a second one makes BibTeX drop every reference.
\documentclass[11pt]{article}
% Not [review]: that option anonymises the author block to "Anonymous ACL
% submission" and adds line numbers. This is a course submission, so the names
% belong on it. Switch back to [review] only for an actual blind submission.
\usepackage{acl}

\usepackage{times}
\usepackage{latexsym}
\usepackage[T1]{fontenc}
\usepackage[utf8]{inputenc}
\usepackage{microtype}
\usepackage{inconsolata}
\usepackage{graphicx}
\usepackage{booktabs}
\usepackage{multirow}
\usepackage{amsmath}

\title{VeriMem: Poisoning the Memory of a Retrieval-Augmented Fact-Checking Agent}

% TODO before submitting: add the other four team members and the real emails.
\author{
  Teesha Ramchandani \\
  Indian Institute of Technology Bhilai \\
  \texttt{<email>}
}

\begin{document}
\maketitle

\begin{abstract}
Retrieval-augmented agents increasingly store their own past trajectories and
retrieve them as context for new inputs. We ask what an adversary gains from
writing to that store. We build VeriMem, a fact-checking agent that reproduces
CRAG \citep{yan2024crag} and ExpRAG \citep{ferraz2026exprag} from scratch and
transfers both to claim verification on AVeriTeC \citep{schlichtkrull2023averitec},
and we attack it at two surfaces: the experience store and the evidence corpus.
Injecting fabricated trajectories raises attack success from $0.220$ to $0.430$
as the poisoned share rises to 20\%, while an agent with no store stays at
$0.160$. Two defences --- evaluator gating and Bayesian source trust --- each
flatten that curve, but \emph{which one does the work inverts with the
adversary}: against fakes citing unknown domains, trust is decisive and gating
adds nothing; against an adversary that cites the domains the agent already
trusts, trust fails completely ($0.460$, identical to undefended) and gating
carries the defence alone. We show why, by measuring what the trust model
actually learns: it does not separate known misinformation domains from the rest
($0.380$ vs.\ $0.375$), and ranks \texttt{twitter.com} above \texttt{nature.com}.
It is a novelty detector, not a credibility detector. We also report that
poisoning the evidence corpus is both more effective (attack success $0.690$)
and less defensible, because fabricated pages are \emph{genuinely relevant} and
the relevance evaluator scores them almost as highly as real evidence. Both
reproductions produce honest negative results, and we locate the mechanism for
each.
\end{abstract}

\section{Introduction}

A retrieval-augmented agent that keeps a memory of its own past work has a new
attack surface: the memory. An adversary who can write to it shapes how future
inputs are judged without touching the model, the retriever, or the evidence
corpus. The records they write need not be malformed --- a poisoned trajectory
can be syntactically perfect, which is precisely what makes it hard to find.

We study this concretely. VeriMem is a fact-checking agent: given a claim it
returns a verdict in \{\textsc{Supported}, \textsc{Refuted}, \textsc{Conflicting
evidence}, \textsc{Not enough evidence}\} together with an evidence trail ---
the sources it used, their relevance grades, the trust it assigns each
publisher, and the past verifications it drew on. It combines two systems we
reproduce from scratch: CRAG's retrieval evaluator and three-way routing, and
ExpRAG's trajectory store feeding retrieval-augmented fine-tuning.

Our contributions are:

\begin{enumerate}
\itemsep0em
\item \textbf{A memory-poisoning study} (\S\ref{sec:c1}) at 0/5/10/20\%
  injection, with two adversaries and two attack surfaces. We find that
  experience retrieval is a real vulnerability (H1), that defences cut attack
  success (H2), and --- the finding we think matters most --- that the two
  defences are \emph{complementary rather than redundant}, each failing against
  the adversary the other stops.
\item \textbf{Evaluator-gated experience retrieval} (C2), which scores a
  retrieved trajectory on whether its own evidence supports its own verdict.
\item \textbf{Bayesian source trust} (C3), and a measurement showing it learns
  something other than what it was designed to learn (\S\ref{sec:trust}).
\item \textbf{The first reproduction of ExpRAG} we are aware of, which does not
  replicate its central claim, with the reason measured (\S\ref{sec:exprag}).
\end{enumerate}

Every number below carries a bootstrap confidence interval over the evaluation
set. This is not decoration: macro-F1 on AVeriTeC dev gives a quarter of its
weight to a 38-instance class on which every system we tested scores near zero,
and our own early readings of this table were wrong because of it.

\section{Related work}

\paragraph{Corrective retrieval.} CRAG \citep{yan2024crag} adds a fine-tuned
evaluator that grades retrieved documents and routes on the result: refine,
rewrite-and-search, or both. Refinement is decompose-then-recompose over
sentence strips. CRAG reports consistent gains across four QA benchmarks.

\paragraph{Experience replay.} ExpRAG \citep{ferraz2026exprag} stores completed
agent trajectories and uses them twice --- as fine-tuning data, and as retrieved
context during supervised fine-tuning, so the model learns to use precedent.

\paragraph{Attacks on retrieval.} Corpus poisoning against RAG is established
\citep{zou2024poisonedrag, zhong2023poisoning}: an adversary who can insert
documents can steer generation. Our source-poisoning condition (\S\ref{sec:source})
is of this family. The memory-poisoning condition is different in kind: the
adversary writes to the agent's record of \emph{its own reasoning}, and the
record that eventually does the damage may be one the agent wrote itself.

\paragraph{Fact verification.} AVeriTeC \citep{schlichtkrull2023averitec}
provides real-world claims with annotated evidence from archived web pages and a
four-way verdict including two uncertainty labels. The uncertainty labels matter
here: they give a model somewhere to retreat to, which turns out to drive our
dominant failure mode (\S\ref{sec:failures}).

\section{VeriMem}

\subsection{Pipeline}

A claim is verified in five steps. (1)~\textbf{Retrieve}: BM25 over a pooled
corpus built from the archived pages AVeriTeC's annotators cited. (2)~\textbf{Grade}:
a fine-tuned Flan-T5-base evaluator scores each passage for relevance in
$[-1,1]$, read as $P(\text{yes})$ against $P(\text{no})$ at the first decoded
position; the claim's grade is the maximum over passages. (3)~\textbf{Route}:
two thresholds, fitted on 150 development claims, select refine, rewrite, or
both. (4)~\textbf{Retrieve experience} (C2): past trajectories, optionally
gated. (5)~\textbf{Verify}: a verdict from either a prompted teacher or a
LoRA-fine-tuned student, after which source trust is updated (C3) and the
trajectory is stored.

Four configuration flags --- \texttt{experience\_retrieval},
\texttt{gate\_experience}, \texttt{source\_trust}, \texttt{verifier} --- control
the contributions. With all of them off the pipeline must reproduce the Phase~1
baseline exactly; a regression test asserts that passing the full Phase~2--3
machinery to a flags-off run produces byte-identical model calls, and separately
that each flag changes behaviour when enabled, so an inert flag cannot pass.

\subsection{C2: evaluator-gated experience retrieval}

A retrieved trajectory is scored by the same evaluator that grades evidence,
but on a different question: \emph{does this trajectory's evidence support the
verdict it reached?} Our first specification asked instead whether the
trajectory was relevant to the incoming claim, and it was defeated --- a
poisoned record is written \emph{about} the target claim, so it is genuinely
relevant and the evaluator scored it highly, correctly. Relevance is not
truthfulness. Scoring evidence against the verdict it is supposed to justify is
a question a fabrication can fail.

\subsection{C3: Bayesian source trust}

Each domain carries a $\mathrm{Beta}(\alpha,\beta)$ posterior with
$\text{trust} = \alpha/(\alpha+\beta)$, updated by whether the verdict a
trajectory reached was correct. Trajectories and evidence strips whose domains
fall below a floor of $0.3$ are dropped.

The prior is $\mathrm{Beta}(1,3)$, not $\mathrm{Beta}(1,1)$. A uniform prior
places an unseen domain at $0.5$, above any usable floor, so an adversary
registering fresh domains walks straight through; a model that only distrusts
domains it has already watched fail cannot defend against a new attacker. Under
$\mathrm{Beta}(1,3)$ a stranger starts at $0.25$ and must earn its way in. This
choice turns out to be the whole of C3's effect (\S\ref{sec:trust}).

\section{Experimental setup}

\paragraph{Data.} AVeriTeC. The official knowledge store was withdrawn during
the project, so we rebuilt a corpus from the Wayback snapshots in the gold
annotations: 764 documents / 11{,}559 passages for dev, 3{,}468 / 51{,}471 for
train. Evidence is pooled across claims rather than kept per claim; a per-claim
pool makes BM25 unable to fail and the baseline meaningless. Fact-checking
domains are blocklisted, and the blocklist resolves archive wrappers to the
underlying publisher --- 50 of 961 dev URLs were fact-checkers inside Wayback
URLs, which would otherwise have passed.

\paragraph{Models.} Teacher and judge: Qwen2.5-7B-Instruct, served locally.
Student: Qwen2.5-3B-Instruct with QLoRA ($r{=}16$, $\alpha{=}32$, 4-bit base).
Evaluator: Flan-T5-base. All runs are single-seed with bootstrap intervals;
see \S\ref{sec:limitations}.

\paragraph{Metrics.} Macro-F1 and accuracy against the four-way gold label. For
the attacks, \emph{attack success rate} is the fraction of targeted claims that
receive the adversary's chosen verdict, reported separately from collateral
degradation --- a defence that cuts attack success by discarding all retrieved
experience has not defended anything.

\section{Reproducing CRAG}
\label{sec:crag}

\begin{table}[t]
\centering\small
\begin{tabular}{lcc}
\toprule
System & Macro-F1 & 95\% CI \\
\midrule
No retrieval & 0.324 & [0.288, 0.360] \\
Plain RAG    & 0.354 & [0.319, 0.388] \\
CRAG         & 0.363 & [0.324, 0.404] \\
\midrule
\multicolumn{3}{l}{\emph{Paired differences}} \\
Plain RAG $-$ no retr. & $+0.029$ & [$-0.016$, $+0.078$] \\
CRAG $-$ plain RAG     & $+0.009$ & [$-0.031$, $+0.051$] \\
CRAG $-$ no retrieval  & $+0.038$ & [$-0.005$, $+0.085$] \\
\bottomrule
\end{tabular}
\caption{AVeriTeC dev, 500 claims. No pairwise difference excludes zero.}
\label{tab:crag}
\end{table}

\paragraph{The evaluator reproduces.} Fine-tuned Flan-T5-base reaches AUC
$0.811$ / accuracy $0.728$ on 751 held-out pairs, against AUC
$0.650$ $[0.611, 0.686]$ / accuracy $0.573$ for an off-the-shelf
\texttt{ms-marco-MiniLM} cross-encoder: $+0.161$ AUC from fine-tuning. The
\emph{reason} is specific. The baseline's accuracy barely clears the
majority-class floor ($0.531$) while its AUC shows real ranking ability, and its
mean score on genuinely relevant passages is $-0.296$: it ranks adequately and
decides in the wrong place. Since CRAG reads two thresholds off this score, a
misplaced boundary routes nearly everything down one branch however well the
model ranks. Fine-tuning moved the boundary, not only the ordering.

\paragraph{Routing reproduces.} Fitted thresholds beat the best fixed branch by
$0.044$ macro-F1, where guessed thresholds beat it by $0.004$.

\paragraph{Refinement does not transfer.} Table~\ref{tab:crag} shows no
significant pairwise difference. The null is two effects cancelling. Refinement
keeps the best five \emph{sentences} where plain RAG passes the best five
\emph{passages}, so the verifier sees 1{,}281 characters of evidence instead of
3{,}750 --- 66\% less --- and it abstains. CRAG answers \textsc{Not enough
evidence} 145 times against a gold count of 35; 31.0\% of its predictions are an
uncertainty label against 14.6\% in gold, while plain RAG sits at 15.2\%,
almost exactly calibrated. Accuracy falls significantly ($-0.054$
$[-0.096, -0.012]$) while the two rare classes gain, and macro-F1 --- weighting
all four equally --- cancels them.

This is a task-transfer limit rather than a failed reproduction. Refinement was
designed for QA, where a short answer span suffices and surrounding context is
noise. In fact verification the verdict depends on how much the evidence
\emph{establishes}, and \textsc{Not enough evidence} is a valid label to retreat
into. Thinning the evidence by two thirds makes that retreat the easy option.

\section{Reproducing ExpRAG}
\label{sec:exprag}

\begin{table}[t]
\centering\small
\begin{tabular}{llcc}
\toprule
Adapter & Retr. & Macro-F1 & Acc. \\
\midrule
none            & --- & 0.243 & 0.263 \\
none            & yes & 0.285 & 0.333 \\
LoRA SFT        & --- & 0.329 & 0.667 \\
retrieval-aug.\ SFT & yes & 0.332 & 0.673 \\
\midrule
\multicolumn{4}{l}{\emph{filtered targets, for comparison}} \\
LoRA SFT        & --- & 0.302 & 0.630 \\
retrieval-aug.\ SFT & yes & 0.300 & 0.633 \\
\bottomrule
\end{tabular}
\caption{AVeriTeC dev, 300 claims, Qwen2.5-3B student.}
\label{tab:exprag}
\end{table}

\paragraph{Fine-tuning on the agent's own trajectories works.} Accuracy rises
$0.263 \rightarrow 0.667$, macro-F1 by $+0.086$ (Table~\ref{tab:exprag}).

\paragraph{The central claim does not replicate.} ExpRAG argues that training
\emph{with} retrieved trajectories teaches the model to use them, so
retrieval-augmented SFT should beat plain SFT by more than retrieval alone buys.
Decomposed against the prompted baseline: retrieval alone $+0.041$, training
alone $+0.086$, both together $+0.088$ --- a super-additivity of $-0.039$.
Paired, retrieval-augmented SFT $-$ plain SFT is $+0.002$ $[-0.021, +0.026]$:
a tight interval, so a measured null rather than an underpowered test.

The prompted-plus-retrieval row is what makes this falsifiable and is the row
easiest to omit; without it, ``retrieval-augmented SFT is the best row'' cannot
be distinguished from retrieval merely being available at inference.

\paragraph{Why.} The mechanism fired --- all 1{,}500 training examples carried
retrieved context --- but the context was empty of signal. Measuring each
retrieved neighbour's BM25 similarity as a fraction of the claim's own
self-score: median $0.203$, with only 21.0\% above $0.30$ and 7.3\% above
$0.50$. In roughly four training examples out of five the two ``relevant past
verifications'' were barely related to the claim, so the model learned to
discount the context. That retrieval still helps the \emph{prompted} student
($+0.041$) fits the same explanation: it was never trained to ignore it.

ExpRAG's premise requires a store dense enough for precedents to exist. AVeriTeC
spans politics, health and economics across many countries and years, and 1{,}500
trajectories over that spread do not contain near-neighbours for most claims.
The result predicts where the method would work: a narrower domain, or a store
one to two orders of magnitude larger.

\paragraph{Training toward the gold label.} ExpRAG trains on trajectories as
produced. Keeping only the agent's \emph{correct} verifications starved the
minority classes --- of 885 usable examples, 5 were \textsc{Conflicting
evidence}, because those are exactly the claims the agent gets wrong. Training
toward the claim's gold label instead keeps every labelled trajectory and
restores the distribution (\textsc{Conflicting} $5 \rightarrow 106$). It gains
$+0.027$ $[+0.004, +0.052]$ plain and $+0.032$ $[+0.009, +0.056]$ with
retrieval, both significant. It did not, however, survive training: both
fine-tuned models score $0.00$ F1 on both uncertainty labels, while the prompted
model manages $0.13$ and $0.14$. Fine-tuning bought accuracy by abandoning the
uncertain classes.

\section{Memory poisoning}
\label{sec:c1}

\subsection{Threat model}

The adversary can write records into the experience store but cannot modify the
model, the retriever, the evidence corpus, or the gold labels. For a target
claim they write trajectories that are (a) close enough to be retrieved beside
it and (b) confidently assert the verdict they want, citing fabricated evidence.
Every injected record is a syntactically valid trajectory: no format check can
find them. Fabricated records paraphrase rather than copy the target claim,
since an identical claim is removed by self-exclusion before it reaches a
prompt.

We use two adversaries. The \textbf{naive} one cites domains from AVeriTeC's
published misinformation list. The \textbf{adaptive} one reads the agent's own
history and cites the five domains it trusts most --- in our run
\texttt{en.wikipedia.org}, \texttt{twitter.com}, \texttt{cdc.gov},
\texttt{who.int}, \texttt{ncbi.nlm.nih.gov}.

\subsection{H1: retrieval is the vulnerability}

\begin{figure}[t]
\centering
\includegraphics[width=\columnwidth]{figures/memory_poisoning.png}
\caption{Attack success against poisoned share, 100 targets. No-memory has no
store to write to.}
\label{fig:memory}
\end{figure}

With experience retrieval on, attack success climbs monotonically from $0.220$
at 0\% to $0.430$ at 20\% poisoning (Figure~\ref{fig:memory}); the intervals at
the two ends do not overlap. An agent with no store sits at $0.160$. Plain
retrieval also costs $0.060$ attack success and $0.040$ macro-F1 at \emph{zero}
poisoning: trajectory accuracy is $0.590$, so clean experience is often wrong
experience.

\subsection{H2, and the inversion}

\begin{table}[t]
\centering\small
\begin{tabular}{lcc}
\toprule
Configuration & Naive & Adaptive \\
\midrule
memory (undefended) & 0.430 & 0.460 \\
$+$ gating (C2)     & 0.320 & \textbf{0.310} \\
$+$ trust (C3)      & \textbf{0.230} & 0.460 \\
$+$ both            & 0.230 & 0.310 \\
\midrule
no memory           & \multicolumn{2}{c}{0.160} \\
\bottomrule
\end{tabular}
\caption{Attack success at 20\% poisoning. Which defence carries H2 inverts with
the adversary; bold marks the one doing the work.}
\label{tab:inversion}
\end{table}

Both defences flatten the curve, so H2 holds. But Table~\ref{tab:inversion}
shows the result we think is the paper's: \emph{which} defence does it inverts
entirely. Against the naive adversary, trust alone holds attack success at
$0.230$ and keeps zero poisoned records, and gating adds nothing. Against the
adaptive adversary, trust is defeated outright --- $0.460$, identical to
undefended, with every poisoned record retained --- and gating carries the
defence alone at $0.310$.

\begin{figure}[t]
\centering
\includegraphics[width=\columnwidth]{figures/adversary_inversion.png}
\caption{The same cells as Table~\ref{tab:inversion} with their intervals. The
reversal between the two defences is wider than the uncertainty on it; the
narrower comparisons between neighbouring bars are not.}
\label{fig:inversion}
\end{figure}

Each defence fails against the adversary the other stops. This is a direct
argument that C2 and C3 are both necessary; reporting either alone, or only one
adversary, would have supported the opposite conclusion. Figure~\ref{fig:inversion}
shows the same cells with their intervals, which is the honest form of the
claim: trust moving from $0.230$ to $0.460$ between adversaries is far outside
the uncertainty, while the gap between gating and trust at a single adversary is
not.

We also note the limit plainly: gated memory under the adaptive attack is
$0.310$, still nearly double no-memory's $0.160$. The defences reduce the cost
of having a memory. They do not remove it.

\subsection{What source trust actually learns}
\label{sec:trust}

\begin{figure}[t]
\centering
\includegraphics[width=\columnwidth]{figures/trust_divergence.png}
\caption{C3 does not separate misinformation domains from the rest, and the
ordering it does learn is not a credibility ordering.}
\label{fig:trust}
\end{figure}

The inversion has a single explanation, and it is uncomfortable. Fitting trust
on all 1{,}500 trajectories gives 954 domains a track record. Mean trust is
$0.380$ for the 20 domains on AVeriTeC's misinformation list and $0.375$ for the
other 934 --- a difference of $0.005$, with the distributions overlapping
entirely (Figure~\ref{fig:trust}). Among domains with at least ten observations,
\texttt{twitter.com} is the most trusted at $0.76$, while \texttt{nature.com}
sits at $0.29$ and \texttt{cbo.gov} at $0.23$.

No credibility model produces that ordering. C3's does, because it is not a
credibility model: it credits every domain in a trajectory according to whether
the \emph{verdict} was correct, so what accumulates is ``citing this domain
coincided with being right''. Claims evidenced by a CBO report are hard; claims
evidenced by a tweet are often easy.

So C3 stopped the naive adversary not by recognising misinformation but because
those domains were \emph{unseen}, and $\mathrm{Beta}(1,3)$ places a stranger
below the floor. It is a \textbf{novelty detector}. The adaptive adversary
defeats it by ceasing to be novel. The two results are one mechanism.

\subsection{Source poisoning}
\label{sec:source}

\begin{figure}[t]
\centering
\includegraphics[width=\columnwidth]{figures/source_poisoning.png}
\caption{Poisoning the evidence corpus, and the reason gating cannot defend it.}
\label{fig:source}
\end{figure}

Writing fabricated \emph{pages} into the evidence pool is both more effective
and less defensible. Attack success rises $0.190 \rightarrow 0.690$ at five
fabricated pages per target, against memory poisoning's $0.430$ ceiling, and
macro-F1 collapses from $0.374$ to $0.178$ --- unlike memory poisoning, this
attack is broadly destructive rather than surgical.

Gating cannot help, and Figure~\ref{fig:source} measures why rather than
asserting it. A fabricated page is written about the target claim in detail: it
\emph{is} relevant, and the evaluator is right to say so. Mean relevance grade
on fabricated evidence is $0.667$ against $0.740$ on real evidence at one page
per target, narrowing to $0.605$ against $0.643$ at five. Gaps of $0.02$--$0.07$
on a $[-1,1]$ scale leave the gate no signal to act on.

Trust reduces the attack ($0.690 \rightarrow 0.440$ at five pages) but against
the adaptive adversary it does nothing at all: $0.620$ against an undefended
$0.630$, with every fabricated strip retained. Against an adversary that
impersonates trusted publishers, VeriMem has no defence at this surface.

\subsection{Propagation}

If a poisoned record makes the agent misjudge claim $A$, and $A$'s verdict is
then stored, does it contaminate a neighbouring claim $B$? The record doing the
harm at that point is not fabricated --- it is the pipeline's own output citing
real publishers, so C2 sees a coherent trajectory and C3 sees ordinary domains.

We paired each of 100 targets with its nearest different claim by BM25
similarity relative to the target's self-score, excluding duplicates, giving 77
pairs. Each neighbour was verified three times against stores differing
\emph{only} in $A$'s verdict: absent, clean, corrupted.

Of 28 targets the attack actually flipped, \textbf{0 propagated}. The channel
existed --- $A$'s trajectory was retrieved in 67 of 77 pairs under both the
clean and corrupted conditions, and 0 of 77 under absent --- so the poisoned
verdict reached $B$'s prompt and $B$'s answer did not move. With 28 trials the
honest statement is that one-hop propagation is below roughly 10\%, not that it
is zero. We attribute the null to the neighbour being recognisably a different
claim at median relative similarity $0.23$; direct memory poisoning works
because the fake paraphrases its target.

\section{Where the errors come from}
\label{sec:failures}

Categorising all 227 errors on dev by predicates over each trajectory's own
record: \textbf{over-abstention is 52.9\%} --- answering \textsc{Not enough
evidence} where the gold label is definite. False confidence is 20.3\% and
polarity flips 20.3\%.

Only 11 of 227 errors (5\%) had no strip graded above $+0.3$. In the other 95\%
the evaluator judged the evidence relevant and the verifier still erred, so the
failure is in verification, not retrieval --- which contradicts the intuitive
diagnosis. A representative case: a claim graded $0.9825$ and routed
\texttt{correct} received \textsc{Not enough evidence}, the verifier reasoning
that the evidence said nothing about the specific person named. Both components
were right about their own input; the passages were on-topic, which is what the
evaluator scores, and the surviving sentences had lost the specific fact, which
is what the verifier needed.

The clearest next change this points at is passing refined strips \emph{alongside}
more of their surrounding passage rather than instead of it.

\section{Limitations}
\label{sec:limitations}

\paragraph{The evaluator has saturated its own labels.} A human hand-check of 95
sampled pairs puts label noise at $0.289$ (reweighted to the dev distribution;
positives are the worst class at $0.360$). A classifier perfect against truth
would score about $0.711$ accuracy against labels that noisy, and our evaluator
scores $0.728$ --- \emph{above} that ceiling, which means its errors correlate
with the labeller's. Accuracy $0.728$ reads as ``agrees 73\% of the time with a
labelling that is itself 29\% wrong''. The cross-encoder comparison narrows
correspondingly: part of the $+0.161$ AUC is a better match to this labelling
rather than to relevance, and the two cannot be separated without clean labels.
The calibration finding stands independently. CRAG's thresholds were fitted
against the same noisy score.

\paragraph{Single seed.} Every reported number carries a bootstrap interval over
the evaluation set, which answers whether a difference survives a different draw
of claims; it does not capture variance from training seeds.

\paragraph{Evaluation split.} A held-out-\emph{topic} split was planned and not
built. Results are on an AVeriTeC dev subset, which is held out from the
trajectories' source split but not by topic.

\paragraph{Scale.} 1{,}500 trajectories, one 3B student, one 7B teacher. The
ExpRAG result in particular is partly a statement about that budget.

\paragraph{Deviations.} Retrieval is BM25 rather than FAISS over embeddings
(deterministic, and shared with passage retrieval, which matters when comparing
runs differing only in a store's contents). The agent is plain Python rather
than a graph framework. Evidence strips were not deduplicated when these runs
were made; repeats cost 5.2\% of evidence slots, measured, and the code was
fixed afterwards without re-running.

\section{Conclusion}

An agent that remembers its own work can be attacked through that memory, and
the attack scales with how much of the store the adversary controls. Two
defences each appear sufficient in isolation and neither is: evaluator gating
and source trust fail against different adversaries, and a study using only one
adversary would have concluded that either alone suffices. The trust model in
particular does not do what it was designed to do --- it detects novelty, not
unreliability --- which is both why it stops an attacker using unknown domains
and why it offers nothing against one that cites the agent's own favourites.

\section*{Ethics statement}

All fabricated pages and trajectories exist only to attack our own system. They
are generated with fixed seeds into a gitignored directory, assembled in memory
for a run and discarded; no fabricated content is written into the evidence
corpus on disk, served, or published. Attack domains are drawn from AVeriTeC's
own published list of known misinformation sites, so the study does not invent
new disinformation and does not attribute fabricated claims to a publisher that
has done nothing. We report a genuine weakness in a class of system that is
being deployed; we judge that describing the attack surface and two partial
defences is more useful than withholding it, particularly as the stronger
adaptive attack requires only reading a system's own public source history.

% acl.sty already issues \bibliographystyle{acl_natbib}; a second one makes BibTeX
% skip the command and silently drop every reference.
\bibliography{verimem}

\end{document}
